\section{Pembahasan}\label{sec:discussion}
\subsection{Memisahkan Manfaat Fine-tuning dari Manfaat Pre-training}
Pada TS2Vec, keunggulan atas heuristik ada baik pada encoder acak maupun pada encoder hasil pre-training di EX-1, sehingga keunggulan tersebut tidak dapat dikaitkan dengan pre-training. Salah satu penjelasan yang masuk akal adalah bahwa encoder konvolusional yang diinisialisasi acak sudah menyediakan fitur yang berguna ketika sebuah head dilatih di atasnya, tetapi penjelasan ini merupakan hipotesis dan belum diuji di sini. \todo{Putuskan apakah akan menambah eksperimen (misalnya linear probe pada fitur mentah) untuk menguji hipotesis ini. Konsekuensi: tanpa uji tersebut, penjelasan hanya boleh ditulis sebagai dugaan.}

\subsection{Kapan Transfer Membantu dan Kapan Tidak}
Satu-satunya efek transfer yang teramati sejauh ini adalah RE1 $\rightarrow$ RE2: efeknya kecil dan tidak terlihat pada 10\% label, dengan RE3 sebagai sumber, maupun dengan RE1 sebagai target. RE1 dan RE2 memiliki tipe fault yang sama kecuali fault socket, sehingga representasi RE1 mungkin sesuai untuk RE2, sedangkan RE3 memuat fault tingkat kode. Penjelasan berbasis kesamaan tipe fault ini adalah hipotesis yang perlu diperiksa, misalnya dengan membandingkan hasil per tipe fault. \todo{Laporkan hasil per tipe fault jika data memungkinkan. Konsekuensi: tanpa prediksi per kasus yang tersimpan, analisis ini membutuhkan eksperimen ulang.} Mode FineAll tidak lebih baik daripada FineLast; laporan kemajuan mengaitkannya dengan overfitting yang cepat, dan hal ini perlu diverifikasi dengan kurva pelatihan.

\subsection{Membaca Perbandingan Antar Metode}
\todo{Setelah eksperimen: (i) laporkan C1 lebih dulu sebagai perbandingan praktis; (ii) laporkan keuntungan setiap metode terhadap encoder acaknya sendiri (C2) dan nyatakan apakah perbedaannya melampaui ketidakpastiannya; (iii) bahas apakah perbedaan berasal dari tugas pretext, backbone, atau anggaran pre-training; (iv) bahas RUN-encoder secara terpisah dari RUN-original dan jangan menyimpulkan apa pun tentang RUN dari RUN-encoder.}
Aturan keputusan yang ditetapkan sebelumnya: pre-training dinyatakan bermanfaat hanya jika paired bootstrap terhadap encoder acak pada mode yang sama signifikan, dan dinyatakan tidak memiliki efek yang terdeteksi jika tidak.

\subsection{Keterbatasan}
Data uji kecil (75, 45, dan 18 kasus), sehingga selisih satu atau dua kasus berada dalam rentang derau. Pembagian data hanya satu dan seed berjumlah tiga. Bootstrap hanya me-resample kasus. Waktu injeksi diasumsikan diketahui. Heuristik dan semua encoder hanya memakai lima channel metrik. RE3 jenuh sehingga tidak informatif. Hasil saat ini hanya untuk TS2Vec, sehingga belum ada kesimpulan tentang kelebihan relatif ketiga metode. Generalisasi lintas dataset masih terbuka karena transfer hanya membantu pada satu arah.
